\section{AddRandomToSystem} \label{sec:AddRandomToSystem}

This utility takes an existing system specified by \tt{vtf} file(s) and
adds new beads into it, placing them either randomly or using provided
coordinates. The utility requires an input coordinate file containing all
beads (i.e., all beads defined in the structure file must be present in the
coordinate). If an incomplete coordinate file is provided,
\tt{AddRandomToSystem} exhibits undefined behaviour, i.e., it can run without
problems, exit with an error, or freeze.

The utility reads what is to be added from a
\tt{FIELD}-like file, placing the new beads and molecules randomly.
The random coordinates of added beads can be constrained via several
options; for added molecules, either the molecule's first bead (default
behaviour) or its geometric centre (\tt{--gc} option) obey the
constraints. The coordinates of the remaining beads in the molecule are
governed by the coordinates in the \tt{FIELD}-like file. Therefore,
not all the molecular beads necessarily obey the constraining rules.
Molecules are added with a random orientation.

There are two types of constraints which can be combined: place new beads
(i) specified distance from other beads or (ii) in a specified interval in
x, y, and/or z directions. In (i), options \tt{-ld} and/or \tt{-hd}
specify the distance; if present, these must be accompanied by \tt{-bt}
option.  Then, the new beads are placed at least \tt{-ld <float>} and
at most \tt{-hd <float>} distance from beads specified by the
\tt{-bt} option. In (ii), options \tt{-cx}, \tt{-cy}, and
\tt{-cz} basically change the box size for the added beads; for
example, \tt{-cx 5 9} would generate x coordinates only in the interval
$\langle5,9)$. All these options can be combined, but note that
\tt{AddRandomToSystem} does not perform any sanity checks, that is, if any
combination of the provided options is impossible to achieve, the utility
will run forever.

As stated above, the structure and number of added molecules and monomeric
beads are read from a \tt{FIELD}-like file. This file must contain a
\tt{species} section followed by a \tt{molecule} section as
described in the DL\_MESO simulation package. Following is a short
description of the two sections.

The new unbonded beads appear in the resulting files after
the unbonded beads from the original system (but before any bonded beads).
If molecules are added, \tt{AddRandomToSystem} places them at the end.

By default, the new beads are added to the system, increasing the total
number of beads. However, if \tt{--switch} is used, the new beads
replace beads in the original system. The replaced beads are either neutral
beads with the lowest indices (as per ordering in the provided structure
file) or beads of the type specified using \tt{-xb} option. In both
cases, the beads to be replaced must be unbonded.

The size of the new simulation box can be changed using the \tt{-b}
option. There are no constraints on the size of the cuboid box; if the new
box is small, not all beads from the original system are necesserily inside
the box (although the added once always are). The origin of the new box
coincides with the origin of the original box unless \tt{--centre}
option is used; in that case, the centres of the two boxes coincide.

\tt{AddRandomToSystem} creates \tt{vcf} and \tt{vsf} files containing
the new system; the coordinates can also be written into an \tt{xyz}
file.

Several examples of using the utility are provided in the
\tt{Examples/AddRandomToSystem} directory.

\vspace{1em}
\noindent
Usage: \tt{AddToSystem <input.vcf> <out.vsf> <out.vcf> <options>}
\noindent
\begin{longtable}{p{0.235\textwidth}p{0.709\textwidth}}
  \toprule
  \multicolumn{2}{l}{Mandatory arguments} \\
  \midrule
  \tt{<input.vcf>} & input coordinate file (either \tt{vcf} or
    \tt{vtf} format) \\
  \tt{<out.vsf>} & output \tt{vsf} structure file for the new
    system \\
  \tt{<out.vcf>} & output \tt{vcf} coordinate file for the new
    system \\
  \toprule
  \multicolumn{2}{l}{Non-standard options} \\
  \midrule
  \tt{-st <int>} & \tt{<input.vcf>} timestep to add new beads
    to (default: 1) \\
  \tt{-xyz <name>} & also save coordinates to \tt{xyz} file \\
  \tt{--switch} & replace original beads instead of increasing the
    total number of beads \\
  \tt{-b <x> <y> <z>} & side lengths of the new simulation box \\
  \tt{--centre} & place the original simulation box in the middle of
    the new one \\
  \midrule
  \multicolumn{2}{c}{Random placement} \\
  \midrule
  \tt{-f <name>} & \tt{FIELD}-like file specifying additions to the
    system (default: \tt{FIELD}) \\
  \tt{-ld <float>} & lowest distance from beads specified by
    \tt{-bt} option \\
  \tt{-hd <float>} & highest distance from beads specified by
    \tt{-bt} option \\
  \tt{-bt <bead names>} & bead types to use in conjunction with
    \tt{-ld} and/or \tt{-hd} options \\
  \tt{-cx <num> <num2>} & constrain x coordinate to interval
    $\langle num, num2)$ \\
  \tt{-cy <num> <num2>} & constrain y coordinate to interval
    $\langle num, num2)$ \\
  \tt{-cz <num> <num2>} & constrain z coordinate to interval
    $\langle num, num2)$ \\
  \tt{--gc} & use molecule's geometric centre for distance check \\
  \midrule
  \multicolumn{2}{c}{Provided coordinates} \\
  \midrule
  \tt{-vtf <vsf> <vcf>} & add ready-made system from \tt{vtf} files
    instead of randomly placing beads read from the \tt{FIELD}-like file \\
  \bottomrule
\end{longtable}
